
\medskip

%On étudie les précipitations (hauteurs de pluies) sur la ville de Nouméa entre avril et décembre 2020. 
%
%On obtient le tableau suivant:
%
%\begin{center}
%\begin{tabularx}{\linewidth}{|m{3.5cm}|*{9}{>{\centering \arraybackslash}X|}}\hline
%Mois 				&Avril 	&Mai &Juin 	&Juillet &Août 	&Sept 	&Oct &Nov 	&Déc.\\ \hline
%Précipitations en mm&147 	&199 &40 	&67 	&47 	&54 	&104 &45 	&63\\ \hline
%\multicolumn{10}{r}{\footnotesize Source: https://www.historigue-meteo.net/oceanie/nouvelle-caledonie/noumea/2020}\\
%\end{tabularx}
%\end{center}

\medskip

\begin{enumerate}
\item %Calculer la moyenne des précipitations. Arrondir le résultat au mm près.
La moyenne des précipitations est $\dfrac{147 + 199 + \ldots + 63}{9} = \dfrac{766}{9} \approx 85,1$, soit 85~mm au mm près.
\item %Quelle est l'étendue des précipitations ?
L'étendue est égale à $199 - 40 = 159$.
\item %Déterminer la médiane des précipitations.
En rangeant dans l'ordre croissant : 40 \: 45 \:47 \: 54 \: 63 \: 67 \: 104 \:147 \: 199 on constate que la 5\up{e} valeur soit 63 est la médiane.
\item %Calculer le pourcentage de mois pour lesquels les précipitations sont supérieures à $100$~mm. Arrondir le résultat à l'unité près.
Il y a 3 valeurs supérieures à 100 sur 9 relevés, soit un pourcentage de $\dfrac{3}{9} \times 100 = \dfrac13 \times 100 = \dfrac{100}{3} \approx 33,333$, soit 33\,\% à l'unité près.
\end{enumerate}

\bigskip

